\section{本章小结}
本文提出了一种新的Cosserat杆表示方法，可以同时完美满足不可伸长约束、单位四元数约束和对齐约束。基于该表示方法，我们在SQP框架下仿真Cosserat杆，并进行了两项技术改进以获得更好的性能：首先，我们指出在链式表示下，频繁的矩阵-向量乘法可以在$\mathcal{O}(n)$时间内实现；其次，我们在$\mathcal{O}(n)$时间内构造了一个特定的混合预条件子。结果表明，与块对角预条件子相比，我们的方法在大时间步长和复杂场景下获得了一到两个数量级的加速，并且精确满足非线性约束。

我们方法的一个局限性是，如\Cref{fig:PCG_iter_Young}所述，混合预条件子在$\J_{\q,\s}^\top \mathbf{H}_\q \J_{\q,\s}$主导$A$时表现良好，而当其影响减小时，混合预条件子的优势会降低，此时需要增加PCG迭代次数才能收敛。我们方法的另一个局限性是我们专注于非环形Cosserat杆结构，我们的方法需要进一步扩展到Cosserat网状结构或多分支结构~\cite{spillmann2008cosseratnet,deul_direct_2018}，我们将此作为本文的未来工作。展望未来，我们还想将自适应离散化~\cite{wen_cosserat_2020,sanchez-banderas_robust_2020}纳入我们的方法，以更好地捕捉具有尖锐接触的复杂行为。{\color{black}此外，我们仅在每次SQP求解开始时使用离散碰撞检测来收集邻近对。这种策略可能会错过时间步之间的碰撞以及内部牛顿步骤之间的碰撞。这可能会带来穿透错误或相交。目前我们使用相对较小的时间步来缓解这个问题。更好的方法可以是采用连续策略~\cite{otaduy2009implicit,Smith_2015_Bijective,Li_2020_incremental,wang2021large,Yu_2021_repulsive_curves}。}最后，我们希望研究更高效的求解器或在GPU上并行化我们的方法以获得更好的性能。